\chapter{Конструкторская часть}

В данном разделе описана архитектура компилятора и основные алгоритмы обработки программы без привязки к деталям конкретной реализации.

\section{Общая архитектура}
Компилятор имеет конвейерную архитектуру.
Исходный текст сначала передается лексеру и парсеру ANTLR4.
Полученное конкретное синтаксическое дерево преобразуется в абстрактное синтаксическое дерево.
Затем выполняется семантический анализ, после которого backend генерирует текст CIL.

\includeimage
	{compiler_components}
	{f}
	{H}
	{\textwidth}
	{Компоненты компилятора Plankalkul}

Основные компоненты системы:
\begin{itemize}
	\item фронтенд, отвечающий за разбор текста и построение AST;
	\item модуль AST, задающий узлы типов, выражений, операторов и программ;
	\item семантический анализатор, формирующий таблицы символов и проверяющий типы;
	\item генератор CIL, преобразующий проверенную программу в текстовую сборку;
	\item управляющий модуль компилятора и интерфейс командной строки.
\end{itemize}

\section{Абстрактное синтаксическое дерево}
AST используется как промежуточное представление программы.
В отличие от конкретного дерева ANTLR, оно не содержит служебных синтаксических узлов и хранит только смысловые элементы языка.

Корневым узлом является программа, содержащая список определений типов и список определений программ.
Определение программы хранит номер \texttt{P}, необязательное имя, список параметров, список результатов и тело.
Тело состоит из операторов присваивания, guarded-операторов, \texttt{W}-циклов и \texttt{FOR}-циклов.
Выражения представлены литералами, ссылками на переменные, вызовами программ, унарными и бинарными операциями.

\section{Семантический анализ}
Семантический анализ выполняется в три фазы.
На первой фазе регистрируются пользовательские типы.
На второй фазе регистрируются сигнатуры всех программ, что позволяет поддерживать вызовы программ, объявленных ниже по тексту.
На третьей фазе проверяются тела программ.

\includeimage
	{semantic_phases}
	{f}
	{H}
	{\textwidth}
	{Фазы семантического анализа}

Для каждой программы формируется область видимости.
Параметры класса \texttt{V} и результаты класса \texttt{R} объявляются по сигнатуре программы.
Промежуточные переменные \texttt{Z} и константы \texttt{C} получают тип при первом использовании с inline-аннотацией или из типа присваиваемого выражения.
Для переменных класса \texttt{C} дополнительно проверяется запрет повторного присваивания.

Типы после разрешения представлены структурно: бит, битовая последовательность, массив и запись.
Совместимость присваивания проверяется по структуре типа.
Для компонентного доступа анализатор последовательно проходит по цепочке компонентов и вычисляет тип выбранного поля записи, элемента массива или бита последовательности.

\section{Генерация CIL}
Каждая программа Plankalkul отображается на статический метод класса \texttt{PlankalkulProgram}.
Параметры класса \texttt{V} становятся аргументами метода.
Переменные \texttt{Z}, \texttt{C} и \texttt{R} размещаются в локальных переменных CIL.
Первый по номеру метод используется как целевая программа для точки входа \texttt{Main}.

Записи отображаются на CLR value type с полями \texttt{f\_0}, \texttt{f\_1} и далее.
Если программа возвращает несколько результатов, генератор создает синтетическую запись и возвращает ее как единое значение.
Массивы создаются инструкцией \texttt{newarr}; доступ к элементам и полям реализуется инструкциями \texttt{ldelem}, \texttt{stelem}, \texttt{ldfld}, \texttt{stfld}.

Арифметические и логические выражения транслируются в стековые инструкции CIL.
Guard-оператор реализуется ветвлением \texttt{brfalse}.
Цикл \texttt{FOR} транслируется в пару меток начала и конца цикла, а \texttt{W}-цикл -- в повторное выполнение тела, пока хотя бы одно условие продолжает выполняться.

\section*{Вывод по конструкторской части}
Была спроектирована модульная архитектура компилятора, разделяющая синтаксический анализ, промежуточное представление, семантическую проверку и генерацию CIL.
Определены основные структуры данных и правила отображения конструкций Plankalkul на элементы CLR.
